\section*{第 \ref{ch:b2:comparison} 章——函数的比较}

\begin{solution}{exo:b2:comparison:1}
$\sqrt{x^2 + x + 1} = x\sqrt{1 + \tfrac1x + \tfrac{1}{x^2}} = x +
\frac12 + \frac38\cdot\frac1x + o\bigl(\frac1x\bigr)$（二项式
展开：$\frac12 u - \frac18 u^2$，其中 $u = \frac1x +
\frac{1}{x^2}$ 得 $\frac{1}{2x} + \frac{1}{2x^2} -
\frac{1}{8x^2} = \frac{1}{2x} + \frac{3}{8x^2}$，再乘以
$x$）。

$\ln(x^2 + x) - 2\ln x = \ln\bigl(1 + \tfrac1x\bigr) = \frac1x -
\frac{1}{2x^2} + o\bigl(\frac{1}{x^2}\bigr)$。

第三个函数：将各因子展开，
\[
\frac{x + \sin x}{x - \ln x}
= \Bigl(1 + \frac{\sin x}{x}\Bigr)
\Bigl(1 + \frac{\ln x}{x} + \frac{(\ln x)^2}{x^2} +
O\Bigl(\frac{(\ln x)^3}{x^3}\Bigr)\Bigr).
\]
在 $+\infty$ 的\omterm{def:b2:comparison:landau}{比较尺度}上将各项排序：$\frac{\ln x}{x}
\gg \frac{1}{x} \geq \bigl|\frac{\sin x}{x}\bigr| \gg \frac{(\ln
x)^2}{x^2}$。于是前导 $1$ 之后的两项为
$\frac{\ln x}{x}$，随后是有界振荡项 $\frac{\sin
x}{x}$：
\[
\frac{x + \sin x}{x - \ln x}
= 1 + \frac{\ln x}{x} + \frac{\sin x}{x}
+ O\Bigl(\frac{(\ln x)^2}{x^2}\Bigr).
\]
\end{solution}

\begin{solution}{exo:b2:comparison:2}
$f(t) = t^\alpha$（$t \geq 1$）。

$\alpha > -1$：发散，并由
\cref{thm:b2:comparison:seriesintegral}~(2)，可得 $\sum_{k\leq n}
k^\alpha = \frac{n^{\alpha+1}}{\alpha+1} + C + o(1)$，当 $\alpha <
0$ 时（此时 $f$ 递减）；对于 $\alpha \geq 0$（$f$ 递增）时，
用相反的不等号作同样的夹逼，得 $\sum_{k \leq n}
k^\alpha \sim \frac{n^{\alpha + 1}}{\alpha + 1}$。

$\alpha = -1$：$H_n = \ln n + \gamma + o(1)$
（见\cref{ex:b2:comparison:harmonic}）。

$\alpha < -1$：收敛，余项为 $\sum_{k > n} k^\alpha
\sim \frac{n^{\alpha+1}}{-(\alpha+1)}$，由（1）中的夹逼得（两个
积分界都等价于该值）。
\end{solution}

\begin{solution}{exo:b2:comparison:3}
设 $v_n = H_n - \ln n - \gamma \to 0$。则
\[
v_n - v_{n+1} = \ln\frac{n+1}{n} - \frac{1}{n+1}
= \Bigl(\frac1n - \frac{1}{2n^2}\Bigr) - \Bigl(\frac1n -
\frac{1}{n^2}\Bigr) + O\Bigl(\frac{1}{n^3}\Bigr)
= \frac{1}{2n^2} + O\Bigl(\frac{1}{n^3}\Bigr),
\]
利用 $\frac{1}{n+1} = \frac1n - \frac{1}{n^2} + O(n^{-3})$。因
$v_n \to 0$，对尾部作裂项求和：
\[
v_n = \sum_{k \geq n} (v_k - v_{k+1})
= \sum_{k\geq n} \Bigl(\frac{1}{2k^2} + O(k^{-3})\Bigr)
= \frac{1}{2n} + O\Bigl(\frac{1}{n^2}\Bigr),
\]
由\cref{thm:b2:comparison:seriesintegral}~(1)，用于 $t^{-2}$
（余项 $\sim \frac1n$，减半）及 $t^{-3}$。因此 $H_n =
\ln n + \gamma + \frac{1}{2n} + o(\frac1n)$。
\end{solution}

\begin{solution}{exo:b2:comparison:4}
三次应用斯特林公式：
\[
\frac{(3n)!}{(n!)^3}
\sim \frac{\sqrt{6\pi n}\,(3n/\eu)^{3n}}
{(2\pi n)^{3/2}\,(n/\eu)^{3n}}
= \frac{\sqrt{6}\; 27^{\,n}}{2\pi n} \cdot
\frac{1}{\sqrt{2\pi n}}\cdot\sqrt{2\pi n}\;
= \frac{\sqrt3\,27^n}{2\pi n} .
\]
（注意：$\frac{\sqrt{6\pi n}}{(2\pi n)^{3/2}} =
\frac{\sqrt6}{(2\pi n)\sqrt{2\pi n}}\sqrt{\pi n} =
\frac{\sqrt3}{2\pi n}$）。

$\dfrac{n!}{n^n} \sim \sqrt{2\pi n}\,\eu^{-n}$。

$\sqrt[n]{n!} = \exp\bigl(\frac{\ln n!}{n}\bigr)$，其中 $\ln n! = n\ln
n - n + \frac12\ln(2\pi n) + o(1)$：
\[
\sqrt[n]{n!} = \exp\Bigl(\ln n - 1 + \frac{\ln(2\pi n)}{2n} +
o\Bigl(\frac{\ln n}{n}\Bigr)\Bigr)
= \frac{n}{\eu}\Bigl(1 + \frac{\ln(2\pi n)}{2n} +
o\Bigl(\frac{\ln n}{n}\Bigr)\Bigr).
\]
\end{solution}

\begin{solution}{exo:b2:comparison:5}
$g(x) = x^n + x - 1$ 在 $\intcc{0}{1}$ 上严格递增，从 $-1$
增至 $1$：唯一根为 $x_n$。因 $x_n^n = 1 - x_n \in
\intoo{0}{1}$：若沿某子列 $x_n \leq c < 1$，则 $x_n^n
\leq c^n \to 0$，于是 $1 - x_n \to 0$，与 $x_n \leq
c$ 矛盾。因此 $x_n \to 1$。

记 $x_n = 1 - \varepsilon_n$，$\varepsilon_n \to 0^+$。方程为
$(1 - \varepsilon_n)^n = \varepsilon_n$，即
\[
n\ln(1 - \varepsilon_n) = \ln \varepsilon_n
\quad\Longrightarrow\quad
-n\varepsilon_n\bigl(1 + o(1)\bigr) = \ln\varepsilon_n .
\]
于是 $n\varepsilon_n = -\ln\varepsilon_n\,(1 + o(1)) \to +\infty$，并
再取对数，得 $\ln n + \ln\varepsilon_n =
\ln(-\ln\varepsilon_n) + o(1)$。因 $\ln(-\ln \varepsilon_n) =
o(\ln(1/\varepsilon_n))$，由此得 $\ln\varepsilon_n \sim -\ln
n$，从而 $\varepsilon_n = \frac{-\ln\varepsilon_n}{n}(1 + o(1))
\sim \frac{\ln n}{n}$：
\[
x_n = 1 - \frac{\ln n}{n} + o\Bigl(\frac{\ln n}{n}\Bigr) .
\]
\end{solution}

\begin{solution}{exo:b2:comparison:6}
精确关系：$\cot y_n = x_n = n\pi + \frac\pi2 - y_n$，其中 $y_n
\sim \frac{1}{n\pi}$（见\cref{ex:b2:comparison:tan}）。展开 $\cot y
= \frac1y - \frac y3 + O(y^3)$:
\[
\frac{1}{y_n} - \frac{y_n}{3} + O(y_n^3) = n\pi + \frac\pi2 - y_n
\quad\Longrightarrow\quad
\frac{1}{y_n} = n\pi + \frac\pi2 + O\Bigl(\frac1n\Bigr),
\]
（其中 $-y_n$ 与 $-\frac{y_n}{3}$ 均为 $O(\frac1n)$）。取倒数：
\[
y_n = \frac{1}{n\pi}\cdot\frac{1}{1 + \frac{1}{2n} + O(n^{-2})}
= \frac{1}{n\pi}\Bigl(1 - \frac{1}{2n} + O\Bigl(\frac{1}{n^2}\Bigr)\Bigr)
= \frac{1}{n\pi} - \frac{1}{2n^2\pi} + O\Bigl(\frac{1}{n^3}\Bigr).
\]
由此
\[
x_n = n\pi + \frac{\pi}{2} - y_n
= n\pi + \frac\pi2 - \frac{1}{n\pi} + \frac{1}{2n^2\pi} +
o\Bigl(\frac{1}{n^2}\Bigr).
\]
\end{solution}

\begin{solution}{exo:b2:comparison:7}
分出最大的两项：
\[
\sum_{k=0}^{n} k! = n! + (n-1)! + \sum_{k \leq n-2} k! ,
\qquad
\sum_{k\leq n-2} k! \leq (n-1)\,(n-2)! = (n-1)! .
\]
故 $1 \leq \frac{1}{n!}\sum k! \leq 1 + \frac{2}{n}$，极限为
$1$。精细化：$\frac{(n-1)!}{n!} = \frac1n$，粗略的界
$\sum_{k \leq n-2}k! \leq (n-1)!$ 可用同样方式加细：
$\sum_{k\leq n-2} k! = (n-2)!\,(1 + O(\frac1n)) = O\bigl(\frac{n!}{n^2}\bigr)$.
由此
\[
\sum_{k=0}^{n} k! = n!\Bigl(1 + \frac1n + O\Bigl(\frac{1}{n^2}\Bigr)\Bigr).
\]
\end{solution}

\begin{solution}{exo:b2:comparison:8}
$(u_n)$ 递增；若它有界，则将收敛到 $\ell$，且
$\ell = \ell + \frac1\ell$，矛盾。故 $u_n \to \infty$。

取平方：$u_{n+1}^2 = u_n^2 + 2 + u_n^{-2}$，于是
\[
u_n^2 = u_0^2 + 2n + \sum_{k=0}^{n-1} \frac{1}{u_k^2} .
\]
和式为 $o(n)$（各项趋于 $0$，由切萨罗平均），故 $u_n^2 \sim 2n$
且 $u_n \sim \sqrt{2n}$。

精细化：$\frac{1}{u_k^2} \sim \frac{1}{2k}$，于是由比较
（见\cref{thm:b2:comparison:seriesintegral}，或
正项级数部分和的等价量）$\sum_{k<n} u_k^{-2} \sim \frac12 \ln n$。
由此
\[
u_n^2 = 2n + \frac{\ln n}{2}\,(1 + o(1)) + O(1)
\quad\Longrightarrow\quad
u_n = \sqrt{2n}\sqrt{1 + \frac{\ln n}{4n} + o\Bigl(\frac{\ln
n}{n}\Bigr)}
= \sqrt{2n}\Bigl(1 + \frac{\ln n}{8n} + o\Bigl(\frac{\ln
n}{n}\Bigr)\Bigr).
\]
\end{solution}

\begin{solution}{exo:b2:comparison:9}
提出 $n$，并令 $t = \ln n$：
\[
S_n = \frac1n \sum_{k=1}^{n} \frac{1}{1 + t\,\frac kn} .
\]
固定 $t$ 时，和式是函数 $u \mapsto \frac{1}{1 +
tu}$ 在 $\intcc{0}{1}$ 上的黎曼和；该函数关于 $u$ 单调，故
其黎曼和被相差一个步长的积分夹住：
\[
\int_0^1 \frac{\dd u}{1 + tu} - \frac1n
\leq S_n \leq \int_0^1 \frac{\dd u}{1 + tu} + \frac1n
\]
（单调函数的黎曼和与其积分的比较，
对每个 $n$ 以各自的 $t = \ln n$ 成立）。现有
$\int_0^1 \frac{\dd u}{1 + tu} = \frac{\ln(1 + t)}{t}$，且
$\frac1n = o\bigl(\frac{\ln t}{t}\bigr)$，因此
\[
S_n = \frac{\ln(1 + \ln n)}{\ln n} + O\Bigl(\frac 1n\Bigr)
\;\sim\; \frac{\ln\ln n}{\ln n} .
\]
\end{solution}

\begin{solution}{exo:b2:comparison:10}
恒等式：$(\ln n)^{\ln n} = \eu^{\ln n\,\ln\ln n} =
\bigl(\eu^{\ln n}\bigr)^{\ln\ln n} = n^{\ln\ln n}$。排序：
比较对数。$\ln(n^2) = 2\ln n$；$\ln\bigl((\ln
n)^{\ln n}\bigr) = \ln n\ln\ln n$; $\ln(2^n) = n\ln2$;
$\ln(n!) = n\ln n - n + O(\ln n)$（斯特林公式，或较粗糙的
夹逼 $\ln n! \sim n\ln n$）；$\ln(n^n) = n\ln n$。由于
$2\ln n = o(\ln n\ln\ln n)$, $\ln n\ln\ln n = o(n)$, $n\ln 2 =
o(n\ln n - n)$，且 $n \ln n - n \sim n\ln n$，但 $n! / n^n \to
0$（对数之差为 $-n + O(\ln n) \to -\infty$）：
\[
n^2 = o\bigl((\ln n)^{\ln n}\bigr),\quad
(\ln n)^{\ln n} = o(2^n),\quad
2^n = o(n!),\quad
n! = o(n^n).
\]
（对每一步：对数之差趋于
$+\infty$，于是比值趋于 $0$）。
\end{solution}

\begin{solution}{exo:b2:comparison:11}
将 $\frac1{k^2} = \frac1{k(k+1)} + \frac1{k^2(k+1)}$ 分解，并
对 $k > n$ 求和：
\[
\sum_{k>n}\frac1{k^2} = \frac1{n+1} +
\sum_{k>n}\frac{1}{k^2(k+1)} ,
\]
第一个和式恰好裂项相消（$\frac1{k(k+1)} = \frac1k -
\frac1{k+1}$）。对第二个和式：$\frac{1}{k^2(k+1)} = \frac1{k^3}
+ O\bigl(\frac1{k^4}\bigr)$（因为 $\frac{1}{k^2(k+1)} -
\frac1{k^3} = \frac{-1}{k^3(k+1)}$），并由
积分比较，$\sum_{k>n}\frac1{k^3} = \frac1{2n^2} +
O\bigl(\frac1{n^3}\bigr)$, $\sum_{k>n}\frac1{k^4} =
O\bigl(\frac1{n^3}\bigr)$。由此
\[
\sum_{k>n}\frac1{k^2}
= \frac1{n+1} + \frac{1}{2n^2} + O\Bigl(\frac1{n^3}\Bigr)
= \frac1n - \frac1{n^2} + \frac{1}{2n^2} +
O\Bigl(\frac1{n^3}\Bigr)
= \frac1n - \frac{1}{2n^2} + O\Bigl(\frac1{n^3}\Bigr),
\]
利用 $\frac1{n+1} = \frac1n - \frac1{n^2} +
O\bigl(\frac1{n^3}\bigr)$.
\end{solution}

\begin{solution}{exo:b2:comparison:12}
$(u_n)$ 单调递增；若它有界，则会收敛于某个
有限的 $\ell$，且 $\ell = \ell + \eu^{-\ell}$，这不可能。故
$u_n \to \infty$。令 $v_n = \eu^{u_n} \to \infty$，则
\[
v_{n+1} = \eu^{u_n + \eu^{-u_n}} = v_n\,\eu^{1/v_n}
= v_n\Bigl(1 + \frac1{v_n} + \frac1{2v_n^2} +
O\bigl(v_n^{-3}\bigr)\Bigr)
= v_n + 1 + \frac{1}{2v_n} + O\bigl(v_n^{-2}\bigr).
\]
先对 $v_{k+1} - v_k = 1 + O(1)$ 求和，得 $v_n = n +
O(n)$，因此最终有 $v_n \geq cn$；重新求和时用
$\frac1{2v_k} = O(\frac1k)$，得 $v_n = n + O(\ln n)$。
再作一轮：$\frac{1}{2v_k} = \frac{1}{2k}\bigl(1 +
O\bigl(\tfrac{\ln k}k\bigr)\bigr)$，于是
\[
v_n = n + \sum_{k<n}\frac1{2k} + O(1) = n + \frac{\ln n}2 +
O(1).
\]
最后，$u_n = \ln v_n = \ln n + \ln\Bigl(1 + \frac{\ln n}{2n} +
O\bigl(\tfrac1n\bigr)\Bigr) = \ln n + \frac{\ln n}{2n} +
O\bigl(\tfrac1n\bigr)$。
\end{solution}

\begin{solution}{pb:b2:comparison:1}
\textbf{1.} 两式相减，得 $\sum_i (c_i -
c_i')\varphi_i = o(\varphi_k)$。若某个系数不同，设
$i_0$ 为第一个这样的下标：除以 $\varphi_{i_0}$，
利用 $\varphi_j = o(\varphi_{i_0})$（$j > i_0$）得 $c_{i_0} -
c_{i_0}' = o(1)$，故为零，矛盾。关于展开式：
令 $u = \frac{\ln x}x \to 0$，
\[
\frac{1}{x - \ln x} = \frac1x\cdot\frac{1}{1 - u}
= \frac1x\bigl(1 + u + u^2 + O(u^3)\bigr)
= \frac1x + \frac{\ln x}{x^2} + \frac{(\ln x)^2}{x^3} +
o\Bigl(\frac{(\ln x)^2}{x^3}\Bigr).
\]
展开式中不出现 $\frac c{x^2}$ 项，因为该展开是
关于 $u = \frac{\ln x}{x}$ 的几何级数：每一项中 $\ln x$ 的幂次，
与 $\frac1x$ 除首个因子外的幂次相同；
\omterm{def:b2:comparison:landau}{比较尺度}上的 $\frac1{x^2}$ 一级（即 $(\ln x)^0$ 的系数）
恰好缺失，其系数为 $0$。

\textbf{2.} $f(x) = \eu^x + x$ \omterm{def:b2:metric:continuity}{连续}、严格
递增，极限为 $-\infty$ 与 $+\infty$，故为
$\R \to \R$ 的双射，因此 $x_n = f^{-1}(n)$ 存在且唯一，并且
$x_n \to +\infty$（$f^{-1}$ 递增至 $+\infty$）。由
$\eu^{x_n} = n - x_n$ 得 $x_n = \ln(n - x_n) \leq \ln n$，故
$x_n/n \to 0$，且 $x_n = \ln n + \ln(1 - x_n/n) = \ln n + o(1)
\sim \ln n$。

\textbf{3.} 记 $u_n = x_n/n$。第二轮：$u_n = \frac{\ln n
+ o(1)}{n}$，于是
\[
x_n = \ln n + \ln(1 - u_n) = \ln n - u_n + O(u_n^2)
= \ln n - \frac{\ln n}{n} + o\Bigl(\frac{\ln n}n\Bigr).
\]
第三轮：此时 $u_n = \frac{\ln n}{n} - \frac{\ln n}{n^2} +
o\bigl(\frac{\ln n}{n^2}\bigr)$，且 $\ln(1 - u_n) = -u_n -
\frac{u_n^2}2 + O(u_n^3)$：
\[
x_n = \ln n - \frac{\ln n}n + \frac{\ln n}{n^2}
- \frac{(\ln n)^2}{2n^2} + o\Bigl(\frac{(\ln
n)^2}{n^2}\Bigr)
= \ln n - \frac{\ln n}{n} - \frac{(\ln n)^2}{2n^2} +
o\Bigl(\frac{(\ln n)^2}{n^2}\Bigr),
\]
其中项 $\frac{\ln n}{n^2}$ 被吸收进
$o\bigl(\frac{(\ln n)^2}{n^2}\bigr)$ 之中。

\textbf{4.} 当 $n = 1000$ 时：$\ln 1000 \approx 6.90776$（误差
$7\cdot10^{-3}$）；取两项：$6.90085$（误差
$2\cdot10^{-5}$）；取三项：$6.90082$（误差小于
$10^{-5}$），而 $x_{1000} \approx 6.90083$。每一轮
误差大致按预测的因子 $\frac{\ln n}{n}$ 缩小。

\textbf{5.} 从右端开始，\omterm{def:b2:metric:continuity}{连续}两次分部积分：
由 $\frac{\dd}{\dd t}\bigl[\tfrac12t(1-t)\bigr] = \tfrac12 -
t$，且 $t(1-t)$ 在两端点处为零，
\[
\frac12\int_0^1 t(1-t)g''(t)\dd t
= -\int_0^1\Bigl(\frac12 - t\Bigr)g'(t)\dd t
= -\Bigl[\Bigl(\frac12 - t\Bigr)g\Bigr]_0^1 - \int_0^1 g
= \frac{g(0) + g(1)}2 - \int_0^1 g .
\]
移项即得所述恒等式。

\textbf{6.} 先计算增量，再将第 5 题应用于
$g(t) = f(n + t)$：
\begin{align*}
E_{n+1} - E_n
&= f(n{+}1) - \int_n^{n+1}\!f - \frac{f(n{+}1) - f(n)}2 \\
&= \frac{f(n) + f(n{+}1)}2 - \int_n^{n+1}\!f
= \frac12\int_0^1 t(1-t)f''(n+t)\dd t .
\end{align*}
由于 $0 \leq t(1-t) \leq \frac14$，有 $\abs{E_{n+1} - E_n} \leq
\frac18\int_n^{n+1}\abs{f''}$，其关于 $n$ 的和由假设收敛，
因此 $(E_n)$ 收敛（增量绝对可和）
到某个 $E$，并且
\[
\abs{E - E_n} \leq \sum_{k\geq n}\abs{E_{k+1} - E_k} \leq
\frac18\int_n^\infty\abs{f''} .
\]

\textbf{7.} 取 $f(t) = \frac1t$，则 $f''(t) = \frac2{t^3}$，
$\int_1^\infty\abs{f''} = 1 < \infty$。由第 6 题得
\[
H_n = \ln n + \frac{1 + \frac1n}{2} + E + (E_n - E)
= \ln n + \Bigl(E + \frac12\Bigr) + \frac1{2n} +
\varepsilon_n,
\]
其中 $\abs{\varepsilon_n} = \abs{E_n - E} \leq
\frac18\int_n^\infty\frac{2\dd t}{t^3} = \frac1{8n^2}$。
与 $H_n = \ln n + \gamma + o(1)$ 比较，
（见\cref{ex:b2:comparison:harmonic}）可定出 $E + \frac12 =
\gamma$。

\textbf{8.} 由第 6 题的增量公式，
\[
\varepsilon_n = E_n - E = -\sum_{k\geq n}\frac12\int_0^1
t(1-t)\,\frac{2\,\dd t}{(k+t)^3}
= -\sum_{k \geq n}\Bigl(\frac1{k^3}\int_0^1t(1-t)\dd t +
O\Bigl(\frac1{k^4}\Bigr)\Bigr),
\]
利用 $\frac{1}{(k+t)^3} = \frac1{k^3} +
O\bigl(\frac1{k^4}\bigr)$ 在 $t \in \intcc01$ 上一致成立。已知
$\int_0^1 t(1-t) = \frac16$ 及
$\sum_{k\geq n}\frac1{k^3} \sim \frac{1}{2n^2}$
（见\cref{thm:b2:comparison:seriesintegral}）：
\[
\varepsilon_n = -\frac16\cdot\frac{1}{2n^2} +
o\Bigl(\frac1{n^2}\Bigr) = -\frac{1}{12n^2} +
o\Bigl(\frac{1}{n^2}\Bigr).
\]

\textbf{9.} 设 $f = \ln$，$f''(t) = -\frac1{t^2}$，绝对
可积。由第 6 题得
\[
\ln n! = \int_1^n\ln t\,\dd t + \frac{\ln n}2 + E + O\Bigl(
\frac1{8}\int_n^\infty\frac{\dd t}{t^2}\Bigr)
= \Bigl(n + \frac12\Bigr)\ln n - n + 1 + E +
O\Bigl(\frac1n\Bigr),
\]
故 $d_n = 1 + E + O\bigl(\frac1n\bigr)$，$(d_n)$ 的收敛性
——\cref{thm:b2:comparison:stirling}的第一步——以及
速率 $O(1/n)$。（由斯特林公式给出的极限值可得 $E =
\ln\sqrt{2\pi} - 1$。）

\textbf{10.} 设 $f(t) = t^{-1/2}$，$f''(t) = \frac34 t^{-5/2}$，
绝对可积。由第 6 题得：
\[
\sum_{k=1}^n \frac1{\sqrt k}
= 2\sqrt n - 2 + \frac{1 + \frac1{\sqrt n}}2 + E +
O\bigl(n^{-3/2}\bigr)
= 2\sqrt n + c + \frac{1}{2\sqrt n} +
O\bigl(n^{-3/2}\bigr),
\]
其中 $c = E - \frac32$。当 $n = 10^4$ 时，$2\sqrt n = 200$，$c
\approx -1.46035$，$\frac1{2\sqrt n} = 0.005$，预测值为
$198.54465$，而直接计算得 $\sum_{k\leq10^4}k^{-1/2} =
198.544645\dots$——三项，七位数字吻合。

\textbf{11.} 函数 $t \mapsto t\ln t$ \omterm{def:b2:metric:continuity}{连续}，且
在 $\intco1\infty$ 上严格递增（导数 $\ln t + 1 \geq 1$），
其值从 $0$ 增至 $+\infty$，故存在唯一的 $x_n$，且 $x_n \to
\infty$（否则 $x_n\ln x_n$ 将保持有界）。
对 $x_n\ln x_n = n$ 取对数，得 $\ln x_n + \ln\ln x_n = \ln n$；
由 $\ln\ln x_n = o(\ln x_n)$，两边除以 $\ln x_n$ 得
$\frac{\ln n}{\ln x_n} \to 1$，即 $\ln x_n \sim \ln n$。

\textbf{12.} 由 $x_n = \frac{n}{\ln x_n}$ 及 $\ln x_n \sim
\ln n$，得 $x_n \sim \frac{n}{\ln n}$。再取一轮近似，$\ln\ln x_n =
\ln\bigl(\ln n\,(1 + o(1))\bigr) = \ln\ln n + o(1)$，故 $\ln
x_n = \ln n - \ln\ln n + o(1)$，于是
\[
x_n = \frac{n}{\ln n - \ln\ln n + o(1)}
= \frac{n}{\ln n}\cdot\frac{1}{1 - \frac{\ln\ln n +
o(1)}{\ln n}}
= \frac{n}{\ln n}\Bigl(1 + \frac{\ln\ln n}{\ln n} +
o\Bigl(\frac{\ln\ln n}{\ln n}\Bigr)\Bigr).
\]

\textbf{13.} 当 $n = 10^6$ 时，$\frac{n}{\ln n} \approx 72\,382$
（偏差 $18\%$），两项给出 $\approx 86\,140$（偏差
$1.9\%$），而真值 $x \approx 87\,848$。每轮
的改进因子仅为 $\frac{\ln\ln n}{\ln n} \approx
\frac{2.63}{13.8} \approx 0.19$。对数尺度的收敛
极其缓慢，令人懊恼——
凡涉及素数之处，皆是常态。

\textbf{14.} 恰有 $n$ 个素数 $\leq p_n$（即
$p_1, \dots, p_n$），故 $\pi(p_n) = n$。素数定理
（在本阶段不加证明地承认；见本科三年级分册）给出 $n = \pi(p_n) \sim
\frac{p_n}{\ln p_n}$，即 $p_n \sim n\ln p_n$。这正是
方程 $x\ln x \approx n$ 的逆向使用。取对数得：
$\ln p_n = \ln n + \ln\ln p_n + o(1)$，又 $\ln\ln p_n =
o(\ln p_n)$ 推出 $\ln p_n \sim \ln n$（同第 11 题）。
代回得：
\[
p_n \sim n\ln p_n = n\,\ln n\,\frac{\ln p_n}{\ln n} \sim n\ln
n .
\]

\textbf{15.} (a) 固定 $\varepsilon > 0$；当 $k$ 充分大时，$(1 -
\varepsilon)k\ln k \leq p_k \leq (1 + \varepsilon)k\ln k$。借助
递增函数 $t\ln t$ 的比较
（仿\cref{thm:b2:comparison:seriesintegral}的夹逼方法），
$\sum_{k\leq n}k\ln k = \int_1^n t\ln t\,\dd t + O(n\ln n) =
\frac{n^2\ln n}2 - \frac{n^2}4 + O(n\ln n) \sim
\frac{n^2\ln n}2$，故 $\sum_{k\leq n}p_k = \frac{n^2\ln
n}{2}(1 + O(\varepsilon) + o(1))$ 对每个 $\varepsilon$ 成立：
$\sum_{k\leq n}p_k \sim \frac{n^2\ln n}2$。(b) 由
素数定理，在不超过 $10^{100}$ 的整数中，
所占比例 $\sim \frac{1}{\ln 10^{100}} = \frac1{230.26\dots}$
为素数：任取一个 $100$ 位整数，它为素数的
概率约为 $\frac1{230}$。

\textbf{16.} 在 $\intoo{n\pi}{n\pi + \frac\pi2}$ 上，$g(x) = \tan
x - \frac1x$ \omterm{def:b2:metric:continuity}{连续}且严格递增（$g' = 1 +
\tan^2x + \frac1{x^2} > 0$），且在左端点处 $g \to -\frac1{n\pi} < 0$，
在右端点处 $g \to +\infty$，故恰有一个
根 $x_n$。由 $\tan z_n = \tan x_n = \frac1{x_n} \to 0$
且 $z_n \in \intoo{0}{\frac\pi2}$，故 $z_n \to 0^+$。

\textbf{17.} 由 $\tan z_n \sim z_n$ 及 $\frac1{x_n} \sim
\frac1{n\pi}$，得 $z_n \sim \frac1{n\pi}$。

\textbf{18.} 记 $z_n = \arctan\frac1{x_n}$，又 $\arctan u = u +
O(u^3)$。由 $z_n = O(\frac1n)$ 得：
\[
\frac1{x_n} = \frac{1}{n\pi}\cdot\frac1{1 + \frac{z_n}{n\pi}}
= \frac1{n\pi} - \frac{z_n}{n^2\pi^2} +
O\Bigl(\frac1{n^4}\Bigr)
= \frac1{n\pi} + O\Bigl(\frac1{n^3}\Bigr),
\]
于是 $z_n = \frac1{n\pi} + O\bigl(\frac1{n^3}\bigr)$：
$\frac{c}{n^2}$ 项的系数为 $0$，因为
$\frac1{x_n}$ 的第一项修正本身的量级
为 $\frac{z_n}{n^2} = O(n^{-3})$。

\textbf{19.} 将 $z_n = \frac1{n\pi} + O(n^{-3})$ 代入
前面的展示式：
\[
\frac{1}{x_n} = \frac{1}{n\pi} - \frac{1}{n^3\pi^3} +
O\Bigl(\frac1{n^5}\Bigr),
\]
于是 $z_n = \arctan\frac1{x_n} = \frac1{x_n} -
\frac{1}{3}\Bigl(\frac1{x_n}\Bigr)^3 + O\Bigl(\frac1{n^5}\Bigr)
= \frac1{n\pi} - \frac{1}{n^3\pi^3} - \frac{1}{3n^3\pi^3} +
O\Bigl(\frac1{n^5}\Bigr)$：
\[
x_n = n\pi + \frac{1}{n\pi} - \frac{4}{3\pi^3n^3} +
O\Bigl(\frac1{n^5}\Bigr).
\]

\textbf{20.} 当 $n = 3$ 时：一项为 $9.53088$，三项为
$9.52929$，真根为 $9.52933$：误差分别为 $1.5\cdot10^{-3}$ 与
$5\cdot10^{-5}$。对比：对 $\tan x = x$，根必须使
$\tan$ 很大，因此根紧贴窗口的\emph{右}端 $n\pi +
\frac\pi2$ 处，在渐近线之前距离约 $\sim\frac1{n\pi}$
处；对 $x\tan x = 1$，根必须使
$\tan$ 很小，因此根恰位于\emph{左}端 $n\pi$ 之后
不远处，距零点约 $\sim\frac1{n\pi}$。两者方法相同，
但地形恰成镜像。

\textbf{21.} $\sin u < u$ 在 $\intoo0\pi$ 上成立，且 $\sin$ 将
$\intoo0\pi$ 映入 $\intoc01 \subseteq \intoo0\pi$。于是一步之后
$u_1 \in \intoc{0}{1}$，随后 $(u_n)$ 递减，
且以 $0$ 为下界，它收敛于 $\sin$ 的一个不动点，
即收敛于 $0$。展开：$\sin u = u(1 - \frac{u^2}6 +
o(u^2))$，于是
\[
\frac{1}{u_{n+1}^2} - \frac1{u_n^2}
= \frac{1}{u_n^2}\Bigl(\bigl(1 - \tfrac{u_n^2}6 +
o(u_n^2)\bigr)^{-2} - 1\Bigr)
= \frac{1}{u_n^2}\Bigl(\frac{u_n^2}{3} + o(u_n^2)\Bigr)
\longrightarrow \frac13 .
\]

\textbf{22.} 由切萨罗（本科一年级分册），增量的
平均值收敛于同一极限：
\[
\frac{1}{n}\cdot\frac{1}{u_n^2}
= \frac1n\Bigl(\frac1{u_0^2} + \sum_{k=0}^{n-1}
\Bigl(\frac1{u_{k+1}^2} - \frac1{u_k^2}\Bigr)\Bigr)
\longrightarrow \frac13 ,
\]
于是 $u_n^2 \sim \frac3n$，又各项皆为正，故 $u_n
\sim \sqrt{3/n}$。

\textbf{23.} 由第 7 题，$\abs{H_n - \ln n - \gamma -
\frac1{2n}} \leq \frac1{8n^2}$。取 $n = 10^6$，此界为
$\frac{1}{8\cdot10^{12}} = 1.25\cdot10^{-13}$：仅计算三项，
即可把百万项的调和和精确到十三位小数，
并附带完全严格的误差证书——这正是带显式余项的
渐近公式的全部意义所在。

\textbf{24.} (a) 真：$\ln u_n - \ln v_n = \ln\frac{u_n}{v_n}
\to 0$，而 $\ln v_n \to +\infty$，故两对数之比
趋于 $1$。(b) 假：$u_n = n + 1 \sim v_n = n$，但
$\eu^{u_n}/\eu^{v_n} = \eu \neq 1$。等价关系容许的是指数中的
可加误差 $o(1)$，而非 $O(1)$。(c) 假：
$f(x) = x + \sin(x^2) \sim g(x) = x$（在 $+\infty$ 处），但 $f'(x)
= 1 + 2x\cos(x^2)$ 无界振荡，而 $g' = 1$：
等价函数的导数之间未必
有任何可比性。

\textbf{25.} \cref{met:b2:comparison:implicit} 中的循环
三次完全相同地运行：定位根，取出粗略
项，再代回以求下一阶——对 $\eu^x + x = n$
（第一部分）；对 $x\ln x = n$（第三部分）；对 $x\tan x = 1$（第
IV 部分）。梯形修正把级数与积分的比较
由“差收敛”提升为显式的
$\frac{f(1) + f(n)}2$ 项，并附带经过证明的 $O(\int_n^\infty
\abs{f''})$ 余项——得到常数与误差界，而不仅是
收敛性。通往素数的桥梁纯粹是求逆：素数定理
说 $\pi(x)\ln x \approx x$，于是 $p_n$，
由 $\pi(p_n) = n$ 定义，满足 $x\ln x = n$ 型方程——
并继承其渐近性质。第 24 题 (a) 中的规则
使得从 $u_n \sim v_n$ 到 $\ln u_n \sim
\ln v_n$ 的每一步推导均合理（第 11、14 题）；(b) 为假，
这正是不对等价关系取指数的原因。要点：欧拉—麦克劳林
公式的一阶形式（第 6 题），以及素数的渐近律
$p_n \sim n\ln n$，即第 $n$ 个素数（第 14 题）。
\end{solution}
